\clearpage
\SourceLeaf{97}%

\begin{center}
\large [\"UBER DAS WESEN UND DIE DEFINITION DER FUNKTIONEN.]
\end{center}

\begin{center}
\textsc{Gauss} an \textsc{Bessel}, 18. Dezember 1811.
\end{center}

\noindent\hspace*{\parindent}$\mathord{.}\,\mathord{.}\,\mathord{.}\,\mathord{.}\,\mathord{.}\,\mathord{.}$ Zunächst würde ich jemanden, der eine neue Funktion in die Analysis einführen will, um eine Erklärung bitten. Soll sie ausschließlich auf reelle Größen (reelle Werte des Funktionsarguments) angewandt werden, sodass er die imaginären Werte des Arguments gleichsam als bloßes Überbein ansieht --- oder stimmt er meinem Grundsatz zu, dass im Reich der Größen die imaginären $a+b\sqrt{-1}=a+bi$ die gleichen Rechte wie die reellen haben müssen? Um praktischen Nutzen geht es hier nicht. Die Analysis ist für mich eine eigenständige Wissenschaft, die durch die Zurücksetzung jener fingierten Größen außerordentlich an Schönheit und Vollständigkeit verlieren und sich ständig gezwungen sehen würde, sonst allgemein gültige Wahrheiten mit höchst lästigen Einschränkungen zu versehen. $\mathord{.}\,\mathord{.}\,\mathord{.}\,\mathord{.}\,\mathord{.}\,\mathord{.}\,\mathord{.}\,\mathord{.}\,\mathord{.}$ Was soll man sich nun unter $\int \varphi x\mathord{.} dx$ für $x=a+bi$ vorstellen? Will man von klaren Begriffen ausgehen, so muss man offenbar annehmen, dass $x$ durch unendlich kleine Zuwächse (jeder von der Form $a+\beta i$) von dem Wert, für den das Integral 0 sein soll, bis zu $x=a+bi$ übergeht und man dann alle $\varphi x\mathord{.} dx$ \textls{summiert}. Damit ist der Sinn vollständig festgelegt. Der Übergang kann jedoch auf unendlich viele Arten geschehen. So wie man sich das \textls{ganze Reich aller reellen Größen} durch eine unendliche Gerade vorstellen kann, so kann man sich das \textls{ganze Reich aller} Größen, reeller und
